\chapter{Numerische Implementierung}

\section{Beweis 1: Dodekaeder-Orbitale}

\begin{verbatim}
import numpy as np
from scipy.special import sph_harm
from scipy.spatial import KDTree

# Dodekaeder-Ecken (normalisiert)
def dodecahedron_vertices():
    phi = (1 + np.sqrt(5)) / 2
    vertices = []
    for i in range(3):
        for j in range(3):
            for k in range(3):
                if (i + j + k) % 2 == 0:
                    continue
                v = np.array([i, j, k]) - 1
                if np.linalg.norm(v) == 0:
                    continue
                vertices.append(v / np.linalg.norm(v))
    return np.array(vertices)

# Fraktale Gewichtung
def psi_D(theta, phi, l=1.7, D=2.704, l_star=5.3e-11):
    r = 1.0  # Einheitskugel
    Y = sph_harm(0, int(l), theta, phi)
    weight = (r / l_star) ** ((D - 3) / 2)
    return Y * weight

# Berechnung der Wahrscheinlichkeitsdichte auf den Dodekaeder-Ecken
vertices = dodecahedron_vertices()
theta = np.arccos(vertices[:, 2])
phi = np.arctan2(vertices[:, 1], vertices[:, 0])
density = np.abs(psi_D(theta, phi)) ** 2

print("Wahrscheinlichkeitsdichte an den Dodekaeder-Ecken:", density)
\end{verbatim}

\section{Beweis 2: Fraktale Regularisierung der QED}

\begin{verbatim}
import numpy as np
from scipy.integrate import quad

def integrand(k, D=2.704, k_star=1/5.3e-11, alpha=1/137, hbar_c=197e-9):
    F_k = 1.0  # Vereinfachung
    return F_k * (k / k_star) ** (D - 3) / k

def self_energy(D=2.704, k_star=1/5.3e-11):
    alpha = 1/137
    hbar_c = 197e-9  # eV * m
    result, _ = quad(lambda k: integrand(k, D, k_star), 1e-10, np.inf)
    return alpha * hbar_c / np.pi * result

print("Selbstenergie (WDBT+):", self_energy(), "eV")
\end{verbatim}

\section{Beweis 3: Konsistenz der Kalibrierungen}

\begin{verbatim}
import numpy as np

# Konstanten
h = 6.626e-34  # J*s
c = 3e8        # m/s
G = 6.674e-11  # m^3/(kg*s^2)
l0 = 1.8e-15   # m
l_star = 5.3e-11  # m
rho0 = 5e-14   # kg/m^3
D = 2.704

# Vakuumenergiedichte
rho_vac = rho0 * (l0 / l_star) ** (3 - D)
print("Vakuumenergiedichte (WDBT+):", rho_vac, "kg/m^3")

# Umrechnung in J/m^3
rho_vac_J = rho_vac * c**2
print("Vakuumenergiedichte (WDBT+) in J/m^3:", rho_vac_J)
print("Beobachtete Vakuumenergiedichte: ~1e-9 J/m^3")
\end{verbatim}